\documentclass[10pt,a4paper]{amsart}
\usepackage[margin=19mm]{geometry}
\usepackage{amsmath,amssymb,mathtools}
\usepackage[scr=rsfs,cal=euler]{mathalfa}
\usepackage{xfrac}
\usepackage[unicode,hidelinks,bookmarks=false]{hyperref}
\newcommand{\ie}{i.e.\ }
\newcommand{\cf}{cf.\ }
\newcommand{\resp}{respectively\ }
\newcommand{\Subsection}{\subsection}
\providecommand{\nobreakdash}{\nobreak\hbox{-}}
\setlength{\emergencystretch}{2em}
\multlinegap0pt
\usepackage{bm}
\usepackage{bbm}
\usepackage{dsfont}
\usepackage{xspace}
\usepackage{cancel}
\usepackage{mathtools}
\usepackage{amsthm}
\makeatletter
\@ifundefined{theorem}{\newtheorem{theorem}{Theorem}}{}
\@ifundefined{lemma}{\newtheorem{lemma}[theorem]{Lemma}}{}
\@ifundefined{proposition}{\newtheorem{proposition}[theorem]{Proposition}}{}
\makeatother
\def\R{{\mathbb R}}
\newcommand\bN{\mathbb{N}}
\newcommand\bR{\mathbb{R}}
\newcommand\cL{\mathcal{L}}
\newcommand\cO{\mathcal{O}}
\newcommand\calL{\mathcal{L}}
\newcommand{\Domm}{\mathcal{O}}
\newcommand{\Sf}{\mathcal{S}_{f}}
\newcommand{\wt}{\widetilde}
\title[Sharp bounds for non-trace class noise and applications to S]{Sharp bounds for non-trace class noise and applications to SPDEs}
\author{Antonio Agresti, Fabian Germ, Mark Veraar}
\date{Source: arXiv:2601.19639v3 (CC BY 4.0)}
\begin{document}
\maketitle

\noindent\emph{Source and licence.} Sharp bounds for non-trace class noise and applications to SPDEs. Version arXiv:2601.19639v3 (CC BY 4.0), distributed under the Creative Commons Attribution 4.0 International licence, as recorded in the arXiv licence metadata for this exact version and shown on the abstract page https://arxiv.org/abs/2601.19639v3. Full rights notice: licenses/arxiv-2601.19639-CC-BY-4.0.txt.\par\smallskip

\noindent\textit{Editorial scope.} Counted unit: the Theorem whose source label is \texttt{thm:generalONB} in the article (source0.tex lines 2558--2571, source label \texttt{thm:generalONB}), delivered together with the article's own complete original proof of it (source0.tex lines 2642--2693, one \texttt{proof} environment of 52 lines). No mathematics is written, repaired or re-derived: statement and proof are the article's own text line for line. The proof is detached from the statement in the source: the article states the result at lines 2558--2571 and prints its proof at lines 2642--2693, and the proof environment's own header names the statement, so the association is the article's own. Required context retained uncounted: thm:MgTmu delta (lines 1634--1649); lem:zetainfty (lines 1783--1797); eq:boundHsqsharp (lines 1786--1790); prop: interpolation (lines 1851--1876); eq:cond1 (lines 1923--1932); eq:cond2 (lines 1923--1932); eq:cond4 (lines 1923--1932); eq:condtheta (lines 1923--1932); eq:cond3 (lines 1935--1939); eq:cond5 (lines 1935--1939); lem:parameternew (lines 2619--2630); eq:parametersnew lemma (lines 2623--2626). Every definition, display and label of those lines is retained unchanged. Omitted from this package and not redistributed: 1--1633, 1650--1782, 1791--1850, 1877--1922, 1933--1934, 1940--2557, 2572--2618, 2627--2641, 2694--3549. Nothing inside a retained excerpt is omitted or altered: every retained body is delivered exactly as the article's lines are written, so no omission receipt of any kind accompanies this package. Citations: the retained excerpts carry no citation command at all, so no bibliography entry of the article is transcribed and no reference is redistributed. Reference adaptation: none was needed and none was made; every reference inside the delivered unit resolves inside it, so no unresolved reference or citation is produced. Printed slips kept literal: no printed word, symbol, hypothesis or number of the counted statement or of its proof was altered. Layout and class: the article's own class is \texttt{amsart} with packages including amssymb, amsmath, mathtools, mathrsfs, amscd, verbatim, stmaryrd, xcolor, a4wide, enumerate, hyperref, todonotes, forest, tabularx; the wrapper is the lane's pinned \texttt{amsart} editorial wrapper at 10pt on A4, which restates the theorem-like environments the retained text needs and adds the article's own macro definitions that the retained text needs (\texttt{\textbackslash Domm}, \texttt{\textbackslash R}, \texttt{\textbackslash Sf}, \texttt{\textbackslash bN}, \texttt{\textbackslash bR}, \texttt{\textbackslash cL}, \texttt{\textbackslash cO}, \texttt{\textbackslash calL}, \texttt{\textbackslash wt}), with extra wrapper packages (bm, bbm, dsfont, xspace, cancel, mathtools). The delivered numbering is the unit's own. Assets: no retained span contains a figure environment, a graphics-inclusion command, a PDF-page inclusion, a table or a TikZ picture; no image file is copied into this package, the package declares no asset dependency and no image file is redistributed with it. Licence: version arXiv:2601.19639v3 of this article is distributed under the Creative Commons Attribution 4.0 International licence, as recorded in the arXiv licence metadata for this exact version and shown on the abstract page https://arxiv.org/abs/2601.19639v3. A full rights notice is delivered at licenses/arxiv-2601.19639-CC-BY-4.0.txt. Characters: every retained excerpt is pure ASCII, so the delivered engine, XeLaTeX with its default Latin Modern text font, reports no missing character.
\begin{theorem}[Sobolev embeddings for Gaussian series -- Weighted sequences]
\label{thm:MgTmu delta}
Let $\Domm$ be an open set in $\R^d$. 
Let $q\in (1, \infty)$, $\eta\in (1, q)$, $\zeta\in [2, \infty]$
and $s\in(0,d)$ be such that 
\begin{align}\label{eq:conditionsmain delta}
\frac{s}{d} + \frac{1}{q}  \geq \frac{1}{\eta} + \frac{1}{2} - \frac{1}{\zeta}\qquad \text{ and }\qquad 
\frac{1}{\eta} - \frac{1}{\zeta}<\frac{1}{2}.
\end{align}
Then for each $g\in L^{\eta}(\mathcal{O})$ and $\mu\in \ell^\zeta(\Sf)$, the operator 
$M_g R_{\mu}:L^2(\mathcal{O})\to \wt{H}^{-s,q}(\mathcal{O})$ is $\gamma$-radonifying
and 
\begin{align}\label{eq:maingammabound}
\|M_g R_{\mu}\|_{\gamma(L^2(\mathcal{O}), \wt{H}^{-s,q}(\mathcal{O}))}\lesssim_{d,s,q,\eta,\zeta} \|\mu\|_{\ell^{\zeta}(\Sf)} \|g\|_{L^{\eta}(\mathcal{O})}.
\end{align} 
\end{theorem}

\begin{lemma}
\label{lem:zetainfty}
Suppose that
\begin{align}
\label{eq:boundHsqsharp}
s\in \Big(\frac{d}{2},d\Big), 
\qquad 2<\eta<q<\infty, \qquad   \frac{s}{d} + \frac{1}{q} \geq \frac{1}{\eta}+\frac{1}{2}. 
\end{align}
Then for each $g\in L^{\eta}(\mathcal{O})$ and $R\in \calL(L^2(\mathcal{O}))$, the operator $M_g R:L^2(\mathcal{O})\to H^{-s,q}(\mathcal{O})$ is $\gamma$-radonifying and 
\begin{equation}
\label{eq:gamma_L2_Leta_endpoint}
\|M_g R\|_{\gamma(L^2(\mathcal{O}), \wt{H}^{-s,q}(\mathcal{O}))}\lesssim_{d,s,q,\eta} \|g\|_{L^\eta(\mathcal{O})} 
\|R\|_{\calL(L^2(\mathcal{O}))}.
\end{equation}
\end{lemma}

\begin{proposition}
    \label{prop: interpolation}
    Let parameters $\eta_0,\eta_1\in (1,\infty)$,  $\zeta_0, \zeta_1\in [2, \infty]$, $s_0,s_1\in \bR$, $p_0,p_1\in (1,\infty)$ and $q_0,q_1\in (1,\infty)$.
    Let $w_0 = (w_0(k))_{k\geq 1}\in \bR_+^\bN$ and $w_1 = (w_1(k))_{k\geq 1}\in\bR_+^\bN$. 
Suppose  that the bilinear operator
    \begin{align}
        R:=\Bigg\{\begin{split}
        L^{\eta_i}(\R^d)\times \ell^{\zeta_i}_{w_i}&\rightarrow \gamma(L^2(\R^d),H^{s_i,q_i}(\R^d)),\\
        (g,\mu)&\mapsto M_gR_\mu,
        \end{split}
        \quad i = 0,1
    \end{align}
    satisfies $\|R\|_{\cL(L^{\eta_i}\times\ell^{\zeta_i}_{w_i},\gamma (L^2,H^{s_i,q_i}))}\leq C_i$ for $i = 0,1$.
    If $\theta\in (0,1)$ and
    \begin{align}
        \frac{1}{p} = \frac{1-\theta}{p_0}+\frac{\theta}{p_1},\quad s = (1-\theta)s_0 + \theta s_1 
        \quad\text{and}\quad
        w(k) = (w_0(k))^{(1-\theta)\frac{\zeta}{\zeta_0}}(w_1(k))^{\theta\frac{\zeta}{\zeta_1}},\quad k\geq 1,
    \end{align}
    for $(p,p_0,p_1)\in \{(\eta,\eta_0, \eta_1),(\zeta,\zeta_0, \zeta_1),(q,q_0, q_1)\}$, then
    \begin{align}
        R:
        L^{\eta}(\R^d)\times \ell^{\zeta}_w\rightarrow \gamma(L^2(\R^d),H^{s,q}(\R^d)),
    \end{align}
    with norm $\|R\|_{\cL(L^{\eta}\times\ell^{\zeta}_w,\gamma (L^2,H^{s,q}))}\leq C_0^{1-\theta}C_1^\theta$.
\end{proposition}

\begin{align}
\label{eq:condtheta}
\theta &= 1-\frac{2}{\zeta},
\\
    \label{eq:cond1} \frac{1}{q_0} &= \frac{1}{1-\theta} \Big(\frac{1}{q} - \frac{\theta}{q_1}\Big),
\\
\label{eq:cond2} \frac{1}{\eta_0} &= \frac{1}{1-\theta} \Big(\frac{1}{\eta} - \frac{\theta}{\eta_1}\Big),
\\
\label{eq:cond4} s_0 &= \frac{s}{1-\theta}-\frac{\theta s_1}{1-\theta},
\end{align}

\begin{align}
\label{eq:cond3} \frac{s_1}{d} &= - \frac{1}{q_1}  + \frac{1}{\eta_1}+\frac12,
\\
\label{eq:cond5} \frac{1}{\eta_0} &= \frac{s_0}{d} + \frac{1}{q_0}.
\end{align}

\begin{lemma}
\label{lem:parameternew}
Let $\Domm$ be an open set in $\R^d$. 
Suppose that
\begin{align}
\label{eq:parametersnew lemma}
s\geq 0, \quad q\in (1, \infty), \quad   \frac{s}{d} + \frac{1}{q} = \frac{1}{\eta}+\frac12, \quad  \eta\in [2, \infty), 
\end{align}
Then for each $g\in L^{\eta}(\mathcal{O})$ and $\mu\in \ell^2$, the operator $M_g R_{\mu}:L^2(\mathcal{O})\to H^{-s,q}(\mathcal{O})$ is $\gamma$-radonifying and 
\[\|M_g R_{\mu}\|_{\gamma(L^2(\mathcal{O}),H^{-s,q}(\mathcal{O}))}\lesssim \|g\|_{L^\eta(\mathcal{O})} \|\mu\|_{\ell^2},\]
where the constant depends only on the parameters $(d,s,q,\eta)$. 
\end{lemma}

\begin{theorem}[Embeddings for Gaussian series with unweighted sequences]
\label{thm:generalONB}
Let $\Domm$ be an open set in $\R^d$.
Let $s>0$, $q\in (1, \infty)$, $\eta\in [2, \infty)$, $\zeta\in [2, \infty)$ be such that
\begin{align}
\label{eq:parametersnew}
\frac{s}{d} + \frac{1}{q} \geq \frac{1}{\eta} +\frac12\qquad \text{and}  \qquad  \frac{1}{\eta} +\frac{1}{\zeta}>\frac{1}{q}.
\end{align}
Then for each $g\in L^{\eta}(\mathcal{O})$ and $\mu \in \ell^\zeta$, the operator $M_g R_{\mu}:L^2(\mathcal{O})\to \wt{H}^{-s,q}(\mathcal{O})$ is $\gamma$-radonifying and 
\[\|M_g R_{\mu}\|_{\gamma(L^2(\mathcal{O}),\wt{H}^{-s,q}(\mathcal{O}))}
\lesssim_{d,s,q,\eta,\zeta}
         \|g\|_{L^\eta(\cO)} \|\mu\|_{\ell^\zeta}  .
\]
\end{theorem}

\begin{proof}[Proof of Theorem \ref{thm:generalONB}]
Note that the case $\zeta = 2$ is covered by Lemma \ref{lem:parameternew}.
To prove the result, we use multilinear complex interpolation in the parameters $(s, q, \eta, \zeta)$ in a similar way as in Theorem \ref{thm:MgTmu delta}, i.e.\ by applying Proposition \ref{prop: interpolation}.
Note that Lemma \ref{lem:parameternew} gives the result for the parameter $(s_0, q_0, \eta_0, \zeta_0)$ with $\zeta_0 = 2$. Moreover, Lemma \ref{lem:zetainfty} gives the result for $(s_1, q_1, \eta_1, \zeta_1)$ with $\zeta_1=\infty$, where the second condition in \eqref{eq:parametersnew} is automatically satisfied due to
\begin{align}
\label{eq:sharp equality s_1}
    \frac{1}{q_1}<\frac{s_1}{d}+\frac{1}{q_1} = \frac{1}{\eta_1} + \frac{1}{2}.
\end{align}
Let $(s, q, \eta, \zeta)$ be as in the theorem statement. As before \eqref{eq:cond1}, \eqref{eq:cond2}, \eqref{eq:cond3}, \eqref{eq:cond4} need to hold, but here, \eqref{eq:cond5} needs to be replaced by 
\begin{align}
\label{eq:cond5new}
\frac{s_0}{d} + \frac{1}{q_0} = \frac{1}{\eta_0} +\frac12.
\end{align}
This again gives a way to write the unknowns $(s_0, q_0, \eta_0)$ and $(s_1, q_1, \eta_1)$ in terms of $(q_1, \eta_1)$.
Since $\zeta_0 = 2$ and $\zeta_1=\infty$, we again have $\theta = 1-\frac{2}{\zeta}$ as in \eqref{eq:condtheta}, where $\zeta\in (2, \infty)$. 
First, we note that due to \eqref{eq:sharp equality s_1} and \eqref{eq:cond5new}, the first condition in \eqref{eq:parametersnew} is satisfied.
It remains to check that under the second condition in \eqref{eq:parametersnew}, there exist $q_0\in [1, \infty)$, $\eta_0\in [2, \infty)$, $s_0\geq 0$, $s_1>0$, for  $(\eta_1, q_1)$ satisfying $2<\eta_1<q_1$, where $(s_0,q_0,\eta_0,\zeta_0)$ satisfy \eqref{eq:parametersnew lemma} and $(s_1,q_1,\eta_1,\zeta_1)$ satisfy \eqref{eq:boundHsqsharp}.

Due to \eqref{eq:cond2}, the conditions $0<\frac{1}{\eta_0}\leq \frac12$ are equivalent to 
\begin{align}
\label{eq:restrqeta2new}
\frac{1}{\theta \eta} - \frac{1}{2\theta}+\frac12\leq \frac{1}{\eta_1}<\frac{1}{\theta \eta}.
\end{align} 
Due to \eqref{eq:cond1}, the conditions $0<\frac{1}{q_0}\leq 1$ are equivalent to 
\begin{align}
\label{eq:restrqeta3new}
\frac{1}{\theta q} - \frac{1}{\theta}+1\leq\frac{1}{q_1}<\frac{1}{\theta q}.
\end{align} 
Due to \eqref{eq:cond5new}, $s_0\geq 0$ is equivalent to $\frac{1}{2} + \frac{1}{\eta_0}- \frac{1}{q_0}\geq 0$. Due to \eqref{eq:cond1} and \eqref{eq:cond2} this can be rewritten as
\begin{align}\label{eq:restrqeta4new}
\frac{1}{\eta_1}\leq \frac{1}{2\theta} - \frac{1}{2} + \frac{1}{\theta \eta} - \frac{1}{\theta q} + \frac{1}{q_1}. 
\end{align}
Finally, from \eqref{eq:cond3}, we see that $s_1>0$ follows from $q_1>\eta_1$. 
The existence of an admissible $\eta_1$ satisfying  $2<\eta_1<q_1$ together with \eqref{eq:restrqeta2new} and \eqref{eq:restrqeta4new} is then equivalent to 
\begin{align}
\label{eq:restreta_1 combined}
\max\Big\{\frac{1}{q_1}, \frac{1}{\theta \eta} - \frac{1}{2\theta}+\frac12\Big\}<\min\Big\{\frac{1}{\theta \eta}, \frac{1}{2\theta} - \frac{1}{2} + \frac{1}{\theta \eta} - \frac{1}{\theta q} + \frac{1}{q_1}, \frac12\Big\}.
\end{align}
For the right-hand side above to be positive for any $q_1>1$, we require that 
\begin{align}
\label{eq:restrqnew}
\frac{1}{2\theta} - \frac{1}{2} + \frac{1}{\theta \eta} - \frac{1}{\theta q}>0,
\end{align}
Hence, if \eqref{eq:restrqnew} holds,
we can find an admissible $q_1$ such that \eqref{eq:restrqeta3new} holds 
if
\begin{align}\label{eq:restrqnew2}
1 - \frac{1}{\theta} +\frac{1}{\theta q} <\min\Big\{\frac12,\frac{1}{\theta \eta}, \frac{1}{\theta q}\Big\}.
\end{align}
One can check that \eqref{eq:restrqnew} is equivalent to the second condition in \eqref{eq:parametersnew}, due to $\theta = 1-\tfrac{2}{\zeta}$. Moreover, 
the first part of the minimum in \eqref{eq:restrqnew2} leads to $\frac{1}{q} - \frac{1}{\zeta} <\frac12$ which is also true due to the second condition in \eqref{eq:parametersnew} and $\eta\geq 2$. The second part of the minimum in \eqref{eq:restrqnew2} leads to $\frac{1}{\eta}  +\frac{2}{\zeta} > \frac{1}{q}$ which also holds due to the second condition in \eqref{eq:parametersnew}. The third part of the minimum does not lead to any restrictions. 
\end{proof}

\end{document}
